In the High-Luminosity LHC (HL-LHC) era, the instantaneous luminosity is expected to increase by approximately a factor of five relative to current LHC operating conditions.
While these conditions will enable the accumulation of datasets more than an order of magnitude larger than those collected in all preceding LHC runs, they will also substantially increase the number of pileup interactions per event.
To mitigate the resulting degradation of event reconstruction performance, both the CMS and ATLAS experiments are incorporating precision timing measurements into their detectors through a combination of new sub-detectors and upgrades to existing systems, including the electromagnetic calorimeter (ECAL).
These developments enable four-dimensional (4D) event reconstruction techniques that provide powerful tools for pileup suppression and enhanced sensitivity to a wide range of physics signatures, including those involving long-lived particles.

The CMS ECAL, through the fast scintillation response of PbWO$_4$ crystals, the characteristics of the front-end electronics, and the readout sampling rate, is capable of achieving excellent intrinsic timing resolution.
However, ECAL timing reconstruction is sensitive to contributions from pileup interactions, particularly out-of-time pileup (OOT-PU), which becomes a dominant limitation as luminosity increases.
As part of the ECAL performance program, a target timing resolution of approximately 30~ps has been established for HL-LHC operation.
In preparation for these conditions, methods to mitigate the impact of OOT-PU on ECAL timing reconstruction have been studied.

In this work, a cross-correlation (CC) timing reconstruction algorithm is presented.
The CC algorithm exploits the out-of-time pulse amplitudes determined by the pileup-aware multi-fit energy reconstruction together with the known characteristic pulse shape of scintillation signals.
By reconstructing the nominal in-time pulse contribution and determining the pulse time via a cross-correlation with the pulse template, the CC algorithm significantly mitigates the effects of OOT-PU and improves the ECAL timing resolution.

Studies of the mean rechit time as a function of bunch-crossing position demonstrate that the CC algorithm substantially reduces both the dependence of the reconstructed time on bunch-crossing position and the energy dependence of the timing measurement.
In direct comparisons using Run~II data, the CC algorithm improves the same readout unit (SRU) timing resolution by approximately 30~ps and the $Z \rightarrow e^{+}e^{-}$ object-level (ZEE) timing resolution by approximately 35~ps relative to the default ratio-based reconstruction.

In Run~III, the CC algorithm was fully integrated into the CMS software framework and a dedicated online timing calibration was deployed.
A dedicated prompt reconstruction using the CC algorithm as the default ECAL timing method demonstrates an improvement of approximately 80~ps in the SRU resolution relative to the ratio-based reconstruction.
During early Run~III data-taking in 2025, the CC algorithm achieved an SRU resolution of approximately 135~ps and a ZEE resolution of approximately 180~ps in prompt reconstruction.
Based on this demonstrated performance, the CC timing reconstruction was adopted as the default ECAL timing algorithm for CMS data taking in 2025.

With continued operation in Run~III, comparable or improved timing performance is expected.
The CC reconstruction provides a robust and scalable framework for ECAL timing in high-pileup environments and establishes a solid foundation for future precision timing applications in the HL-LHC era.
